\section{Theory: what closure guarantees, and why it can prefer the wrong algebra}\label{sec:theory}
Notation: $\mathfrak a=\spn(A)$, $G=B^\top B$, $P_0$ the orthogonal projector onto $\mathfrak a$, $N=I-P_0$.
Proofs are in Appendix~\ref{app:proofs}. Standard Lie-theoretic facts are cited as such; our contributions are
their consequences for this learning objective.

\paragraph{Invariances.}
With $\varepsilon{=}0$ and $M\in GL(K)$, $\Lc(BM)=K^{-2}\|R(M\otimes M)\|_F^2$, where $R$ is the residual matrix,
so $\sigma_{\min}(M)^4\Lc\le\Lc(BM)\le\sigma_{\max}(M)^4\Lc$. The loss is invariant only under orthogonal
reparameterizations; its zero set (``$\mathfrak a$ is a subalgebra'') is $GL(K)$-invariant. For $\varepsilon>0$,
$P_\varepsilon(BM)=B(G+\varepsilon M^{-\top}M^{-1})^{-1}B^\top$. Because $\operatorname{diag}G=2$, the Gram target
affects no gradient.

\begin{proposition}[Closure, escape and integration]\label{prop:basic}
(a)~If $A_1,\dots,A_K$ are independent and $[A_i,A_j]\in\mathfrak a$, then $\mathfrak a$ is a Lie subalgebra and
integrates to a connected immersed subgroup of $SO(d)$ (standard).
(b)~The uniform escape constant
$\delta=\sup_{X,Y\in\mathfrak a}\|N[X,Y]\|_F/(\|X\|_F\|Y\|_F)$ satisfies
$\delta\le K\sqrt{\Lc(A;0)}/\lambda_{\min}(G)$.
(c)~For $X,Y\in\mathfrak a$ with $\|X\|_F,\|Y\|_F\le r<\frac{\ln2}{4}$,
$\dist(\log(e^Xe^Y),\mathfrak a)\le\frac{\delta}{4}[-\ln(2-e^{4r})-4r]=\tfrac12\delta r^2+O(\delta r^3)$.
(d)~Chained orthogonal transport accumulates per-step errors at most additively.
\end{proposition}

Part (c) is local: it covers rotations of up to $\approx0.12$\,rad, smaller than typical inferred actions.

\paragraph{Closed subalgebras in the benchmarks.}
Write $\so(4)=\su(2)_L\oplus\su(2)_R$ (two commuting ideals).

\begin{proposition}[Classification; standard facts applied to our benchmarks]\label{prop:classify}
(a)~Every 2-dimensional subalgebra of $\so(5)$ is abelian and $SO(5)$-conjugate to the true $\mathfrak t^2$.
(b)~Every 3-dimensional subalgebra of $\so(4)$ is isomorphic to $\su(2)$. It is $SO(4)$-conjugate to exactly one
of three: the diagonal $\so(3)$ acting on $\mathbb R^3\oplus\mathbb R$ (the truth), $\su(2)_L$ or $\su(2)_R$.
Under $O(4)$ the two chiral factors merge into one class.
(c)~No closed 4-dimensional span in $\so(4)$ contains the diagonal $\so(3)$.
\end{proposition}

The three SO(3) classes are separated by the Pfaffian of unit elements ($0$, $+\frac12$, $-\frac12$) and by the
participation ratio ($2$ vs.\ $4$). Diagonal and chiral spans are at subspace distance exactly $\frac12$ (principal
angles of $45^\circ$). An optimization that does not use the classification confirms it: minimizing a span-only
closure residual from \RFlowN{} random 3-dimensional subspaces of $\so(4)$ converged to \RFlowDiag{} diagonal and
\RFlowChiral{} chiral spans, and to nothing else.

\paragraph{Endpoint ambiguity: why the objective can favour a wrong closed algebra.}
BracketNet infers each action from the two endpoint codes, $a=q(z,z')$. For the true SO(3) action this is
ill-posed. A rotation about the axis through $z$ fixes $z$, so $g$ and $g\exp(t\hat z)$ have the same endpoints. A
chiral $\su(2)$ has no such ambiguity: left multiplication by unit quaternions acts simply transitively on every
sphere of $\mathbb R^4\cong\mathbb H$. Call an endpoint rule $q$ \emph{exact} on a triple $(z_0,z_1,z_2)$ if
$\rho(q(z_i,z_j))z_i=z_j$ for $ij\in\{01,12,02\}$ and $\rho(q(z_0,z_2))=\rho(q(z_1,z_2))\rho(q(z_0,z_1))$.

\begin{proposition}[Endpoint ambiguity]\label{prop:endpoint}
On the SO(3) benchmark, use the ground-truth state as the latent code, with an exact decoder.
(a)~With chiral generators, $q(z,z')=\log(z'z^{-1})$ is exact on every triple on a sphere $\|z\|=\text{const}>0$ (all data triples); it is smooth wherever $z'\neq-z$.
The closure, Gram and covariance terms vanish too. Exactness fails only where a coefficient exceeds the bound.
(b)~With the true generators, for every $\epsilon>0$ and every orbit sphere $\{(v,z_4):\|v\|=c\}$, $c>0$, no
continuous rule is exact on all triples in that sphere with pairwise angles $<\epsilon$.
(c)~For any action-inference rule based on endpoints, $\mathbb E\|\rho(q(z,gz))-g\|^2>0$ under the true action:
the stabilizer component of $g$ is not identifiable from $(z,gz)$.
\end{proposition}

Part (b) follows from the hairy-ball theorem: exactness on small triangles would define a continuous section of
$SO(3)\to S^2$. It concerns \emph{exact} zero for \emph{continuous} rules only. A discontinuous rule can be exact
off a null set, and we do not bound how close continuous rules can come (App.~\ref{app:proofs}, Remark). On the data,
ignoring the coefficient bound, the minimal-rotation rule for the true algebra transports exactly but leaves a
composition residual of $3.2\times10^{-3}$; the chiral rule leaves ${\sim}10^{-31}$.

\begin{corollary}[Agreement without correctness]\label{cor:agree}
Two models with identical ground-truth embeddings and $\su(2)_L$ generators have linear CKA $1$ with each other and
with the true state, zero closure, transport and composition loss (up to coefficient clipping), and cross-model
algebra distance $0$, but distance $\frac12$ to the true algebra and inferred actions that are not the true rotations.
\end{corollary}

\paragraph{The implication chain, made precise.}
Write (i) closure, (ii) correct conjugacy class, (iii) recovery after latent alignment, (iv) correct action
inference. (i)$\Rightarrow$(ii) holds for $T^2$ (Prop.~\ref{prop:classify}a) and fails for SO(3)
(Prop.~\ref{prop:classify}b). (ii)$\Rightarrow$(iii) also needs an encoder close to an isometry. (iv) is unattainable
for SO(3) by every endpoint method, even with the true algebra (Prop.~\ref{prop:endpoint}c). Agreement between models
implies none of (ii)--(iv) (Cor.~\ref{cor:agree}).
Proposition~\ref{prop:endpoint} describes the loss landscape at one embedding, and qualitatively. It does not show
that the minimum over trainable rules is lower for chiral solutions, nor which solution gradient descent reaches;
Sec.~\ref{sec:mechanism} tests the latter empirically.
